%\chapter{Longitudinal and Torsional Vibrations of Bars}
%\label{longitudinal-and-torsional-vibrations-of-bars}
149. The next system to the string in order of simplicity is the bar, by
which term is usually understood in Acoustics a mass of matter of
uniform substance and elongated cylindrical form. At the ends the
cylinder is cut off by planes perpendicular to the generating lines. The
centres of inertia of the transverse sections lie on a straight line
which is called the axis.

The vibrations of a bar are of three kinds: longitudinal, torsional, and
lateral. Of these the last are the most important, but at the same time
the most difficult in theory. They are considered by themselves in the
next chapter, and will only be referred to here so far as is necessary
for comparison and contrast with the other two kinds of vibrations.

Longitudinal vibrations are those in which the axis remains unmoved,
while the transverse sections vibrate to and fro in the direction
perpendicular to their planes. The moving power is the resistance
offered by the rod to extension or compression.

One peculiarity of this class of vibrations is at once evident. Since
the force necessary to produce a given extension in a bar is
proportional to the area of the section, while the mass to be moved is
also in the same proportion, it follows that for a bar of given length
and material the periodic times and the modes of vibration are
independent of the area and of the form of the transverse section. A
similar law obtains, as we shall presently see, in the case of torsional vibrations.

It is otherwise when the vibrations are lateral. The periodic
times are indeed independent of the thickness of the bar in the
direction perpendicular to the plane of flexure, but the motive power
in this case, viz. the resistance to bending, increases more rapidly
than the thickness in that plane, and therefore an increase in thickness
is accompanied by a rise of pitch.

In the case of longitudinal and lateral vibrations, the mechanical
constants concerned are the density of the material and the value of
Young's modulus. For small extensions (or compressions) Hooke's law,
according to which the tension varies as the extension, actual length -
natural length
be called e, we have 7'=ge, where g is Young's modulus, and T is the
tension per unit area necessary to produce the extension e. Young's
modulus may therefore be defined as the force which would have to be
applied to a bar of unit section, in order to double its length, if
Hooke's law continued to hold good for so great extensions ; its
dimensions are accordingly those of a force divided by an area.

holds good. If the extension, viz.

The torsional vibrations depend also on a second elastic constant 4,
whose interpretation will be considered in the proper place.

Although in theory the three classes of vibrations, depending
respectively on resistance to extension, to torsion, and to flexure are
quite distinct, and independent of one another so long as the squares of
the strains may be neglected, yet in actual experiments with bars which
are neither uniform in material nor accurately cylindrical in figure it
is often found impossible to excite longitudinal or torsional vibrations
without the accompaniment of some measure of lateral motion. In bars of
ordinary dimensions the gravest lateral motion is far graver than the
gravest longitudinal or torsional motion, and consequently it will
generally happen that the principal tone of either of the latter kinds
agrees more or less perfectly in pitch with some overtone of the former
kind. Under such circumstances the regular modes of vibrations become
unstable, and a small irregularity may produce a great effect. The
difficulty of exciting purely longitudinal vibrations in a bar is
similar to that of getting a string to vibrate in one plane.

With this explanation we may proceed to consider the three classes of
vibrations independently, commencing with longitudinal vibrations, which
will in fact raise no mathematical questions beyond those already
disposed of in the previous chapters.

150. When a rod is stretched by a force parallel to its length, the
stretching is in general accompanied by lateral contraction in such a
manner that the augmentation of volume is less than if the displacement
of every particle were parallel to the axis. In the case of a short rod
and of a particle situated near the cylindrical boundary, this lateral
motion would be comparable in magnitude with the longitudinal motion,
and could not be overlooked without risk of considerable error. But
where a rod, whose length is great in proportion to the linear
dimensions of its section, is subject to a stretching of one sign
throughout, the longitudinal motion accumulates, and thus in the case of
ordinary rods vibrating longitudinally in the graver modes, the inertia
of the lateral motion may be neglected. Moreover we shall see later how
a correction may be introduced, if necessary.

Let 2 be the distance of the layer of particles composing any section
from the equilibrium position of one end, when the rod is unstretched,
either by permanent tension or as the result of vibrations, and let

(1) be the displacement, so that the actual position is given by

(2) The equilibrium and actual position
of a neighbouring layer being w+8e, e+8e++ a dz respectively, the
elongation is a and thus, if 7' be the tension per
unit area acting across the section,

Consider now the forces acting on the slice bounded by 2 and 2+ da. If
the area of the section be o, the tension at a is

by (1) go eo acting in the negative direction, and at x+6z
the tension is
acting in the positive direction; and thus the force on the slice due to
the action of the adjoining parts is on the whole

goo bx.

The mass of the element is pw 62, if p be the original density, and
therefore if X be the accelerating force acting on it, the equation of
equilibrium is X+ ag 0 (2)

In what follows we shall not require to consider the operation of an
impressed force. To find the equation of motion we have only to replace
X by the reaction against acceleration -, and thus if g : p=a', we have

ee (8)

This equation is of the same form as that applicable to the transverse
displacements of a stretched string, and indicates the undisturbed
propagation of waves of any type in the positive and negative
directions. The velocity is relative to the unstretched condition of the
bar; the apparent velocity with which a disturbance is propagated in
space will be greater in the ratio of the stretched and unstretched
lengths of any portion of the bar. The distinction is material only in
the case of permanent tension.

151. For the actual magnitude of the velocity of propagation,

we have ;

which is the ratio of the whole tension necessary (according to Hooke's law)
to double the length of the bar and the longitudinal density. If
the same bar were stretched with total tension 7, and were flexible, the
velocity of propagation of waves along it would be /(Z': pw). In order
then that the velocity might be the same in the two cases, 7' must be
gq, or, in other words, the tension would have to be that theoretically
necessary in order to double the length. The tones of longitudinally
vibrating rods are thus very high in comparison with those obtainable
from strings of comparable length.

In the case of steel the value of g is about 22 x 10 grammes weight per
square centimetre. To express this in absolute units of force on the
c.G. 8.' system, we must multiply by 980. In the same system the density
of steel (identical with its specific gravity referred to water) is 78.
Hence for steel

a= /980 a x 10 = 530,000
approximately, which shews that the velocity of sound in steel is about
530,000 centimetres per second, or about 16 times greater than the
velocity of sound in air. In glass the velocity 4 is about the same as
in steel.

Dangling\footnote{ Centimetre, Gramme, Second. This system is recommended by a Committee
of the British Association. Brit. Ass. Report, 1873.}

It ought to be mentioned that in strictness the value of g determined by
statical experiments is not that which ought to be used here. As in the
case of gases, which will be treated in a subsequent chapter, the rapid
alterations of state concerned in the propagation of sound are attended
with thermal effects, one result of which is to increase the effective
value of qg beyond that obtained from observations on extension
conducted at a constant temperature. But the data are not precise enough
to make this correction of any consequence in the case of solids.

152. The solution of the general equation for the longitudinal
vibrations of an unlimited bar, namely f=f(x-at)+F ( at), being the same
as that applicable to a string, need not be further considered here.
When both ends of a bar are free, there is of course no permanent
tension, and at the ends themselves there is no temporary tension. The
condition for a free end is therefore

To determine the normal modes of vibration, we must assume that
(1) varies as a harmonic function of the time-constant. Then as a function of
, must satisfy da + ne =0 (2),
of which the complete integral is E=A (na t+ B (nw(3), where A and
B are independent of a.

Now since a vanishes always when x =0, we get B=0; and

again since a vanishes when 2= -the natural length of the bar,
(nl) =0, which shews that n is of the form

Pere mercer

t being integral.


Accordingly, the normal modes are given by equations of the form ,

= A (ine (at) se (5),

in which of course an arbitrary constant may be added to c, if desired,

The complete solution for a bar with both ends free is therefore
expressed by

(8),

where A, and B, are arbitrary constants, which may be determined, in the
usual manner, when the initial values of and are given,

A zero value of c is admissible; it gives a term representing a
displacement constant with respect both to space and time, and amounting
in fact only to an alteration of the origin.

The period of the gravest component in (6) corresponding to t=1, is
22--.a, which is the time occupied by a disturbance in travelling twice
the length of the rod. The other tones found by ascribing integral
values to 7 form a complete harmonic scale ; so that according to -this
theory the note given by a rod in longitudinal vibration would be in all
cases musical.

In the gravest mode the centre of the rod, where a= l, is a place of no
motion, or node; but the periodic elongation de
compression de is there a maximum.

153. The case of a bar with one end free and the other fixed may be
deduced from the general solution for a bar with both ends free, and of
twice the length. For whatever may be the initial state of the bar free
at 2=0 and fixed at a=/J, such displacements and velocities may always
be ascribed to the sections of a bar extending from 0 to 27 and free at
both ends as shall make the motions of the parts from 0 to / identical
in the two cases. It is only necessary to suppose that from / to 2/ the
displacements and velocities are: initially equal and opposite to those
found in the portion from 0 to J at an equal distance from the centre
=1, Under these circumstances the centre must by the symmetry remain at
rest throughout the motion, and then the portion from 0 to J satisfies
all the required conditions. We conclude that the vibrations of a bar
free at one end and fixed at the other are identical with those of one
half of a bar of twice the length of which both ends are free, the
latter vibrating only in the uneven modes, obtained by making c in
succession all odd integers. The tones of the bar still belong to a
harmonic scale, but the even tones (octave, etc. of the fundamental) are
wanting.

The period of the gravest tone is the time occupied by a pulse in
travelling four times the length of the bar.

154. Both extremities fixed.  When both ends of a bar are fixed, the conditions
to be satisfied at the ends are that the value of is to be invariable. At =90,
we may suppose that F= . At w=/l, is a small constant a, which is zero if there
be no permanent tension. Independently of the vibrations we have evidently
=a+J,
and we should obtain our result most simply by assuming this term at once, But
it may be instructive to proceed by the general method.

Assuming that as a function of the time varies as

A (n)at + B (n at, we see that as a function of x it must satisfy Es
wE=0, of which the general solution is g= C (nx-+D (nx (L).

But since vanishes with a for all values of c, C=0, and thus we may
write

\[=X \sin(na) +  A \cos(n)at + B e in nat.\]

The condition at 2=1 now gives

\[\sin(nl) +  A \cos(nat) +B \sin(nat) =a,\]
from which it follows that for every finite admissible value of n

(nl) =0, or n= But for the zero value of n, we get

A, (nl) =a,
and the corresponding term in is

\[f= A \sin(ne) =a\]

The complete value of is

\[bat +E r 3 a +B \sin(H) (2),\]

The series of tones form a complete harmonic scale (from which however
any of the members may be missing in any actual case of vibration), and
the period of the gravest component is the time taken by a pulse to
travel twice the length of the rod, the same therefore as if both ends
were free. It must be observed that we have here to do with the
unstretched length of the rod, and that the period for a given natural
length is independent of the permanent tension.

The solution of the problem of the doubly fixed bar in the case of no
permanent tension might also be derived from that of a doubly free bar
by mere differentiation with respect to x. For in the latter problem d /
dx satisfies the necessary differential equation, viz.

\[d \xi/ dx d \xi/ dx d \xi/ dx\]
inasmuch as \(\xi\) satisfies

\[d \xi/ dt =  +/-  a  d \xi/ dx\]
and at both ends \(d \xi /dx\) vanishes. Accordingly \(d \xi /dx\) in
this problem satisfies all the conditions prescribed for in the case
when both ends are fixed. The two series of tones are thus identical.

155. The effect of a small load M attached to any point of the rod is
readily calculated approximately, as it is sufficient to assume the
type of vibration to be unaltered (Article 88). We will take the case of
a rod fixed at =0, and free at c=. The kinetic energy is proportional to

or to

Since the potential energy is unaltered, we see by the principles of
Chapter IV., that the effect of the small load M at a distance a from
the fixed end is to increase the period of the component tones in the
ratio

The small quantity M : l is the ratio of the load to the whole mass of
the rod.

If the load be attached at the free end, (ae) =1, and the
effect is to depress the pitch of every tone by the same small interval.
It will be remembered that c is here an uneven integer.

If the point of attachment of M be a node of any component, the pitch of
that component remains unaltered by the addition.

156.  Another problem worth notice occurs when the load at the free end
is great in comparison with the mass of the rod. In this case we may
assume as the type of vibration, a condition of uniform extension along
the length of the rod.

If be the displacement of the load M, the kinetic energy is

T=4Mi+48 pot dx = 28 (M+ pol)(1).

The tension corresponding to the displacement is gw f , and thus the
potential energy of the displacement is ae (2). The equation of motion
is (M+ a E=0, and if 0 (pe) tr +M (3).

The correction due to the inertia of the rod is thus equivalent to the
addition to M of one-third of the mass of the rod.

\section{Correction for lateral motion} \label{sect:156a}
156 a.
So long as a rod or a wire is uniform, waves of longitudinal
vibration are propagated along it without change of type, but any
interruption, or alteration of mechanical properties, will in general
give rise to reflection. If two uniform wires be joined,
the problem of determining the reflection at the junction may be
conducted as in § 148 a. The conditions to be satisfied at the junction
are (i) the continuity of and (ii) the continuity of
qfad jdxy
measuring
the tension. If pi, p,, 6 i, oi,, Oi, a, denote the volume densities,
the sec .ions, and the velocities in the two wires, the ratio of the
reflected to the incident amplitude is given by

H pitOiai + p io a

The reflection vanishes, or the incident wave is propagated through the
junction without loss, if

pitOiOi = poO aj.... (2).

This result illustrates the difficulty which is met with in obtaining
effective transmission of sound from air to metal, or from metal to air,
in the mechanical telephone. Thus the value of pa is about 100,000 times
greater in the case of steel than in the case of air.
157. Our mathematical discussion of longitudinal vibrations may close
with an estimate of the error involved in neglecting the lateral motion
of the parts of the rod not situated on the
axis. If the ratio of lateral contraction to longitudinal extension be
denoted by y, the lateral displacement of a particle distant r from the
axis will be
wre,
in the case of equilibrium, where e is the extension.
Although in strictness this relation will be modified by the inertia of
the lateral motion, yet for the present, purpose it may be supposed to
hold good.


The constant is a numerical quantity, lying between 0 and 4. If w were
negative, a longitudinal tension would produce a lateral swelling, and
if were greater than 4, the lateral contraction would be great enough to
overbalance the elongation, and cause a diminution of volume on the
whole. The latter state of things would be inconsistent with stability,
and the former can scarcely be possible in ordinary solids. At one time
it was supposed that was necessarily equal to 4, so that there was only
one independent elastic constant, but experiments have since shewn that
y is variable. For glass and brass Wertheim found experimentally

If n denote the lateral displacement of the particle distant r from the
axis, and if the section be circular, the kinetic energy due to the
lateral motion is

Thus the whole kinetic energy is

In the case of a bar free at both ends, we have

Ex (2) BE oe TF in 78 0 d and thus , 2 +30: T=14H 5 .

The effect of the inertia of the lateral motion is therefore to increase
the period in the ratio Cer 4 FF This correction will be nearly
insensible for the graver modes of bars of ordinary proportions of
length to thickness.

[A more complete solution of the problem of the present section has been given
by Pochhammer who applies the general equations for an elastic solid to the
case of an infinitely extended cylinder of circular section. The result for
longitudinal vibrations, so far as the term in is in agreement with that above
determined. A similar investigation has also been published by Chree who has
also treated the more general question in which the cylindrical section is not
restricted to be circular.]

158. Experiments on longitudinal vibrations may be made with rods of
deal or of glass. The vibrations are excited by friction, with a-wet
cloth in the case of glass; but for metal or wooden rods it is necessary
to use leather charged with powdered rosin. The longitudinal vibrations
of a pianoforte string may be excited by gently rubbing it
longitudinally with a piece of India rubber, and those of a violin
string by placing the bow obliquely across the string, and moving it
along the string longitudinally, keeping the same point of the bow upon
the string. The note is unpleasantly shrill in both cases.''

If the peg of the violin be turned go as to alter the pitch of the
lateral vibrations very considerably, it will be found that the pitch of
the longitudinal vibrations has altered very slightly. The reason of
this is that in the case of the lateral vibrations the change of
velocity of wave-transmission depends chiefly on the change of tension,
which is considerable. But in the case of the longitudinal vibrations,
the change of velocity of wave-transmission depends upon the change of
extension, which is comparatively slight'.''

In Savart's experiments on longitudinal vibrations, a peculiar sound,
called by him a \emph{son rauque}, was occasionally observed, whose
pitch was an octave below that of the longitudinal vibration. According
to Terquem the cause of this sound is a transverse vibration, whose
appearance is due to an approximate agreement between its own period and
that of the sub-octave of the longitudinal vibration. If this view be
correct, the phenomenon would be one of the second order, probably
referable to the fact that longitudinal compression of a bar tends to
produce curvature.

159. The second class of vibrations, called torsional, which depend on
the resistance opposed to twisting, is of very small importance. A solid
or hollow cylindrical rod of circular section may be twisted by suitable
forces, applied at the ends, in such a manner that each transverse
section remains in its own plane. But if the section be not circular,
the effect of a twist is of a more complicated character, the twist
being necessarily attended by a warping of the layers of matter
originally composing the normal sections. Although the effects of the
warping might probably be determined in any particular case if it were worth
while, we shall confine ourselves here to the case of a circular section, when
there is no motion parallel to the axis of the rod.

Dangling\footnote{ Donkin's Acoustics, p. 154. 3 Ann. de Chimie, tvu. 129-190,}

The force with which twisting is resisted depends upon an elastic
constant different from q, called the rigidity. If we denote it by n,
the relation between q, n, and may be written

shewing that n lies between g and g. In the case of w=, n=fg.

Let us now suppose that we have to do with a rod in the form of a thin
tube of radius c and thickness dr, and let 9 denote the angular
displacement of any section, distant from the origin.

The rate of twist at 2 is represented by , and the shear of the material
composing the pipe by re. The opposing force per
unit of area is nr and since the area is 2rr dr, the moment
round the axis is

Thus the force of restitution acting on the slice dv has the moment

Now the moment of inertia of the slice under consideration is , and
therefore the equation of motion assumes the form

a6 0 Pp aw. "dat se (Dye

Since this is independent of r, the same equation applies to a cylinder
of finite thickness or to one solid throughout,

Dangling\footnote{ Thomson and Tait, \S 683. This, it should be remarked, applies to
isotropic material only.}

The velocity of wave propagation is / and the whole theory is precisely
similar to that of longitudinal vibrations, the condition
for a free end being af =0, and for a fixed end 6=0, or, if a permanent
twist be contemplated, = constant.

The velocity of longitudinal vibrations is to that of torsional
vibrations in the ratio /q:/n or /(2+2y):1. The same ratio applies to
the frequencies of vibration for bars of equal length vibrating in
corresponding modes under corresponding terminal conditions. If 4 =, the
ratio of frequencies would be

vq: Vn=78 : 3 =168, corresponding to an interval rather greater than a
fifth,

In any case the ratio of frequencies must lie between 2:1=1414, and
/3:1=1782,

1. Longitudinal and torsional vibrations were first investigated by
Chladni.

160. In the present chapter we shall consider the lateral vibrations of
thin elastic rods, which in their natural condition are straight. Next
to those of strings, this class of vibrations is perhaps the most
amenable to theoretical and experimental treatment. There is difficulty
sufficient to bring into prominence some important points connected with
the general theory, which the familiarity of the reader with circular
functions may lead him to pass over too lightly in the application to
strings; while at the same time the difficulties of analysis are not
such as to engross attention which should be devoted to general
mathematical and physical principles.

Daniel Bernoulli\footnote{Comment. Acad. Petrop. t. XIII} seems to have
been the first who attacked the problem. Euler, Riccati, Poisson,
Cauchy, and more recently. Strehlke\footnote{Pogg. Ann, Bd. XXVII},
Lissajous\footnote{Ann. d. Chimie (3), XXX. 385}, and A. Seebeck
\footnote{Abhandlungen d. Math. Phys. Classe d. K. Sachs. Gesellschaft
d. Wissenschaften. Leipzig, 1852} are foremost among those who have
advanced our knowledge of it.

161. The problem divides itself into two parts, according to the
presence, or absence, of a permanent longitudinal tension. The
consideration of permanent tension entails additional complication, and
is of interest only in its application to stretched strings, whose
stiffness, though small, cannot be neglected altogether. Our attention
will therefore be given principally to the two extreme cases, (1) when
there is no permanent tension, (2) when the tension is the chief agent
in the vibration.

With respect to the section of the rod, we shall suppose that one
principal axis lies in the plane of vibration, so that the bending at
every part takes place in a direction of maximum or minimum
(or stationary) flexural rigidity. For example, the surface of the rod may
be one of revolution, each section being circular, though not
necessarily of constant radius. Under these circumstances the potential
energy of the bending for each element of length is proportional to the
square of the curvature multiplied by a quantity depending on the
material of the rod, and on the moment of inertia of the transverse
section about an axis through its centre of inertia perpendicular to the
plane of bending. If w be the area of the section, \(w \kappa^2\) its
moment of inertia, q Young's modulus, ds the element of length, and dV
the corresponding potential energy for a curvature 1/R of the axis of
the rod,
\[ dV=\frac{1}{2}(1/2) q w^{2} ds/R^2 \]
(1) This result is readily obtained by considering the extension of the
various filaments of which the bar may be supposed to be made up. Let 7
be the distance from the axis of the projection on the plane of bending
of a filament of section dw. Then the length of the filament is altered
by the bending in the ratio
\[ 1:1 + R, \]
R being the radius of curvature. Thus on the aide of the axis for which
is positive, viz. on the outward side, a filament is extended, while on
the other side of the axis there is compression. The force necessary to
produce the extension / R is q dw / R by the definition of Young's
modulus; and thus the whole couple by which the bending is resisted
amounts to
\[ q dw / R = w^{2}q/R, \]
if w be the area of the section and its radius of gyration about a line
through the axis, and perpendicular to the plane of bending. The angle
of bending corresponding to a length of axis ds is ds/R, and thus the
work required to bend ds to curvature 1/R is
\[ dV=\frac{1}{2}(1/2) q w^{2} ds/R^2 \]
since the mean is half the final value of the couple.


[For a more complete discussion of the legitimacy of the
foregoing method of calculation the reader must be referred to
works upon the Theory of Elasticity. The question of lateral
vibrations has been specially treated by Pochhammer on the
basis of the general equations.]


For a circular section is one-half the radius.

That the potential energy of the bending would be proportional,
\textit{coeteris paribus}, to the square of the curvature, is evident
beforehand.  If we call the coefficient B, we may take
\[ V= (1/2) B ds/R^2, \]
or, in view of the approximate straightness,
\[ V= (1/2) B (d2 y/ dx2)2 dx \]
(2), in which y is the lateral displacement of that point on the axis of the
rod whose abscissa, measured parallel to the undisturbed position, is In
the case of a rod whose sections are similar and similarly situated B
is a constant, and may be removed from under the integral sign.

The kinetic energy of the moving rod is derived partly from the motion of
translation, parallel to y, of the elements composing it, and partly from the
rotation of the same elements about axes through their centres of inertia
perpendicular to the plane of vibration. The former part is expressed by

(3)

if denote the volume-density. To express the latter part, we have only
to observe that the angular displacement of the element dx is

dy/dx, and therefore its angular velocity (d/dt) dy/dx. The square of
this quantity must be multiplied by half the moment of inertia of the
element, that is, by pw dx. We thus obtain
\[ T = (1/2) dy dx (4).  \]


162. In order to form the equation of motion we may avail ourselves of
the principle of virtual velocities. If for simplicity we confine
ourselves to the case of uniform section, we have

where the terms free from the integral sign are to be taken between the
limits. This expression includes only the internal forces due to the
bending. In what follows we shall suppose that there are no forces
acting from without, or rather none that do work upon the system. A
force of constraint, such as that necessary to hold any point of the bar
at rest, need not be regarded, as it does no work and therefore cannot
appear in the equation of virtual velocities,

(2).

Thus the variational equation of motion is

\(\zeta  de\) (3),

in which the terms free from the integral sign are to be taken between
the limits, From this we derive as the equation to be satisfied at all
points of the length of the bar

while at each end

or, if we introduce the value of B, viz. , and write g + p=',

(4), and for each end

Py .(dy (Py af 545 (Q) + + let B53) y= 0 (B).

Tn these equations 3 expresses the velocity of transmission of
longitudinal waves,

162. - TERMINAL CONDITIONS, - 205

The condition (5) to be satisfied at the ends assumes different forms
according to the circumstances of the case. It is possible to conceive a
constraint of such a nature that the ratio 8 (32) : Sy has a prescribed
finite value. The second boundary condition is then
obtained from (5) by introduction of this ratio. But in all the cases
that we shall have to consider, there is either no constraint
or the constraint is such that either 8 (2) or dy vanishes, and then the
boundary conditions take the form

Py . (dy dy dy 543 () 0, laps - 58) by =0 (6).

We must now distinguish the special cases that may arise. If an end be
free, dy and S 2 are both arbitrary, and the conditions become a
df ds

ay OY pty T=" aide -BU=5 tend (7),

the first of which may be regarded as expressing that no couple acts at
the free end, and the second that no force acts,

If the direction at the end be free, but the end itself be constrained
to remain at rest by the action of an applied force of the necessary
magnitude, in which case for want of a better word the rod is said to be
supported, the conditions are ,

=0, BY =0 v (8),

by which (5) is satisfied. A third case arises when an extremity is
constrained to maintain its direction by an applied couple of the
necessary magnitude, but is free to take any position. We have then

dy Py dy 8 (4) 0, ig - BGA O (9).

- Fourthly, the extremity may be constrained both as to position and
direction, in which case the rod is said to be clamped. The conditions
are plainly

8 (2) = 0, By HO (10).

Of these four cases the first and last are the more important ; the
third we shall omit to consider, as there are no experimental means by
which the contemplated constraint could be realized. Even with this
simplification a considerable variety of problems remain for discussion,
as either end of the bar may be free, clamped or supported, but the
complication thence arising is not so great as might have been expected.
We shall find that different cases may be treated together, and that the
solution for one case may sometimes be derived immediately from that of another.

In experimenting on the vibrations of bars, the condition for a clamped
end may be realized with the aid of a vice of massive construction. In
the case of a free end there is of course no difficulty so far as the
end itself is concerned ; but, when both ends are free, a question
arises as to how the weight of the bar is to be supported. In order to
interfere with the vibration as little as possible, the supports must be
confined to the neighbourhood of the nodal points. It is sometimes
sufficient merely to lay the bar on bridges, or to pass a loop of string
round the bar and draw it tight by screws attached to its ends. For more
exact purposes it would perhaps be preferable to carry the weight of the
bar on a pin traversing hole drilled through the middle of the thickness
in the plane of vibration.

When an end is to be `supported,' it may be pressed into contact with a
fixed plate whose plane is perpendicular to the length of the bar.

163. Before proceeding further we shall introduce a supposition, which
will greatly simplify the analysis, without seriously interfering with
the value of the solution. We shall assume that the terms depending on
the angular motion of the sections of the bar may be neglected, which
amounts to supposing the inertia of each section concentrated at its
centre. We shall afterwards (Article 186) investigate a correction for
the rotatory inertia, and shall prove that under ordinary circumstances
it is small. The equation of motion now becomes
a da GET OO NH (L),
and the boundary conditions for a free end d : T= 7279 (2)

The next step in conformity with the general plan will be the assumption
of the harmonic form of y. We may conveniently take

where J is the length of the bar, and m is an abstract number, whose
value has to be determined. Substituting in ), we obtain

If u=e be a solution, we see that p is one of the fourth roots of unity,
viz +1, -1, +, -; 80 that the complete solution is
u=A (ms)+B (m) +06 + De,
containing four arbitrary constants.

[The simplest case occurs when the motion is strictly periodic
with respect to x, C and D vanishing. If X be the wave-length
and T the period of the vibration, we have
Stt m 2Tr , m x=T'
sO that
T X 2wKb ,(46).]


We have still to satisfy the four boundary conditions, two for each end.
These determine the ratios A: B: 0: D, and furnish besides an equation
which m must satisfy. Thus a series of particular values of m are alone
admissible, and for each m the corresponding u is determined in
everything except a constant multiplier. We shall distinguish the
different functions u belonging to the same system by suffixes,

The value of y at any time may be expanded in a series of the functions
u' (Article 92, 93). If c,, c, etc. be the nominal' coordinates, we have
\(y= \pi e\) (5),
\[ T=  G u + dy y + dx = pw \]
(6). We are fully justified in asserting at this stage that each
integrated product of the functions vanishes, and therefore the process
of the following section need not be regarded as more than a
verification. It is however required in order to determine the value of
the integrated squares.

164. Let u, u' denote two of the normal functions corresponding
respectively to m and m'. Then

or, if dashes indicate differentiation with respect to mo

If we subtract equations (2) after multiplying them by w, u,,
respectively, and then integrate over the length of the bar, we have

(3)

the integrated terms being taken between the limits.

Now whether the end in question be clamped, supported, or free', each
term vanishes on account of one or other of its factors, We may
therefore conclude that, if u, us refer to two' modes of vibration
(corresponding of course to the same terminal conditions) of which a rod
is capable, then

provided m and m' be different.


Dangling\footnote{ The reader should observe that the cases here specified are
particular, and that the right-hand member of (8) vanishes, provided
that au, Bu, tn Gee wt aes gy. Tem hi , Prin 5 dx dz da These conditions
include, for instance, the case of a rod whose end is urged
towards its position of equilibrium by a force proportional to the
displacement, as by a spring without inertia.}

and

The attentive reader will perceive that in the process just followed, we
have in fact retraced the steps by which the fundamental differential
equation was itself proved in Article 162. It is the
original variational equation that has the most immediate connection
with the conjugate property. If we denote y by u and dy
by 2, and the equation in question is:
B as a d+ po = 0(5).
Suppose now
that u' relates to a normal component vibration, so that

ripe = BGS a.

By similar reasoning, if v be a normal function, and u represent any
displacement possible to the system,

now

We conclude that if u and v be both normal functions, which have
different periods,

uv de=0 (6);

and this proof is evidently as direct and general as could be desired.

The reader may investigate the formula corresponding to (6), when the
term representing the rotatory inertia is retained.

By means of (6) we may verify that the admissible values of n are real.
For if n were complex, and u=a+c were a normal function, then a-78, the
conjugate of u, would be a normal function also, corresponding to the
conjugate of n , and then the product of the two functions, being a sum
of squares, would not vanish, when integrated '.

If in (3) m and m' be the same, the equation becomes identically true,
and we cannot at once infer the value of /u,"dz.

Dangling\footnote{ This method is, I believe, due to Poisson.}

We must take m' equal to m + m, and trace the limiting form of
the equation as dm tends to vanish.
[It should be observed that the function is not a normal function of the system;
it is supposed to be derived from by variation of m in (4a) § 163,
the coefficients A,B,C,D being retained constant.]
In this way we find

y 2 Fu du du du ddu dud du

the right-hand side being taken between the limits,

du m u, du x, Now ia etc., im=1'' etc., and thus 2 3m" ul'' Ue ma eae Mg
MD ng 2m, En tue tar (wu ie oe u yeu

in which wv'' = 4, so that between the limits. Now whether an end be clamped,
supported, or free, u au'' = 0, vu" = - 0,
and thus, if we take the origin of at one end of the rod,

tis - x 2 Oa og ju de = F (we - Qui ' + 4 Fl (ul - Que'' +0) +(8).

The form of our integral is independent of the terminal condition at
x=0. If the end e=1 be free, w' and u'\,'' vanish, and accordingly

tae m4 U (1) (9),
that is to say, for a rod with one end free the mean value of u is
one-fourth of the terminal value, and that whether the other end be
clamped, supported, or free.


Again, if we suppose that the rod is clamped at 2 =/, u and w' vanish,
and (8) gives .

t u dx = 31 w"(D

Since this must hold good whatever be the terminal condition at
the other end, we see that for a rod, one end of which is fixed and the
other free,

ih da = 418 (free end) = 1u (fixed end),
shewing that in this case u at the free end is the same as w'\,'' at the
clamped end.

The annexed table gives the values of four times the mean of u in the
different cases,

clamped, free u (free end), or '' (clamped end) free, free u (free end)
clamped, clamped ... u'' (clamped end)

supported, supported - 2 u' w'' (supported end) = 2u'' supported, free u
(free end), or - 2u wx (supported end) supported, clamped '' (clamped
end), or - 2 u u'\,'' (supported end)

By the introduction of these values the expression for T assumes a
simpler form. In the case, for example, of a clamped free or a free-free
rod,

(10), where the end = is supposed to be free,

165. A similar method may be applied to investigate the values of u da,
and I u da. In the derivation of equation (7) of the
preceding section nothing was assumed beyond the truth of the equation
w''\,'' =u, and since this equation is equally true of any of the
derived functions, we are at liberty to replace u by wu' or w'', Thus

mz ul

ak u' da = B u' u + wt 27 i wu i i a

nou ul

14-2

taken between the limits, since the term x'' vanishes in all three
cases,

For a free-free rod

(1),

for, as we shall see, the values of
uu'
must be equal and opposite at the two ends.
Whether u be positive or negative at x = l, u' is positive.
For a rod which is clamped at x=0 and free at x=1 ; +f 0 u dx = 3
(wu'),+ , + (u"u'\,''),.

We have already seen that u,'' = u,, and it will appear (Article 173)
that

Uy'' =-4u,, 80 that

(2), a result that we shall have occasion to use later. By applying the
same equation to the evaluation of f udz, we find

since u's" and wu'' vanish,

Comparing this with (8) Article 164, we see that

= (3),

whatever the terminal conditions may be.

The same result may be arrived at more directly by integrating by parts
the equation


166. We may now form the expression for V in terms of the normal
coordinates.

y= =P lng Tate Gat wf de

(1).

If the functions u be those proper to a rod free at 2 =1, this
expression reduces to

In any case the equations of motion are of the form
\[ \tan() + Be 0 mn, fu do = B, (B)y \]
and, since ,d, is by definition the work done by the impressed forces
during the displacement ,,

= Y (4),

if Y p w dz be the lateral force acting on the element of mass p w da,
If there be no impressed forces, the equation reduces to

+ Dem Py =0 (5),

as we know it ought to do.

167. The significance of the reduction of the integrals u de to
dependence on the terminal values of the function and
its derivatives may be placed in a clearer light by the following line
of argument. To fix the ideas, consider the case of a rod clamped at =0,
and free at =J, vibrating in the normal mode expressed by u. If a small
addition Al be made to the rod at the free end, the form of wu
(considered as a function of az) is changed, but, in accordance with the
general principle established in Chapter IV. (Article 88), we may
calculate the period under the altered circumstances without allowance
for the change of type, if we are content to neglect the square of the
change. In consequence of the straightness of the rod at the place where
the addition is made, there is no alteration in the potential energy,
and therefore the alteration of period depends entirely on the variation
of T. This quantity is increased in the ratio
or 1 : 1 + 1
which is also the ratio in which the square of the period is augmented.
Now, as we shall see presently, the - actual period varies as , and
therefore the change in the square of the period
is in the ratio 1 : 1 +

A comparison of the two ratios shews that

u : u dx = 4 : l.

The above reasoning is not insisted upon as a demonstration, but it
serves at least to explain the reduction of which the integral is
susceptible. Other cases in which such integrals occur may be treated in
a similar manner, but it would often require care to predict with
certainty what amount of discontinuity in the varied type might be
admitted without passing out of the range of the principle on which the
argument depends. The reader may, if he pleases, examine the case of a
string in the middle of which a small piece is interpolated.

168. In treating problems relating to vibrations the usual course has been to
determine in the first place the forms of the normal functions, viz. the
functions representing the normal types, and afterwards to investigate the
integral formulae by means of which the particular solutions may be combined to
suit arbitrary initial circumstances. I have preferred to follow a different
order, the better to bring out the generality of the method, which does not
depend upon a knowledge of the normal functions. In pursuance of the same plan,
I shall now investigate the connection of the arbitrary constants with the
initial circumstances, and solve one or two problems analogous to those treated
under the head of Strings.

The general value of y may be written

y= (4 (: m,t+ B, (mt) u,

+ (4 (m,'t + B, (mt) u, A see (1), so that initially Yo = Ayu, + (2),

\[Yo =  m PB, + MO B y u y +\]

(3).

If we multiply (2) by u, and integrate over the length of the rod, we
get

and similarly from (3)

F = mB. C (5),

formulae which determine the arbitrary constants 4,, B,.

It must be observed that we do not need to prove analytically the
possibility of the expansion expressed by (1). If all the particular
solutions are included, (1) necessarily represents the most general
vibration possible, and may therefore be adapted to represent any
admissible initial state.

Let us now suppose that the rod is originally at rest, in its position
of equilibrium, and is set in motion by a blow which imparts velocity to
a small portion of it. Initially, that is, at the moment when the rod
becomes free, y,=0, and y, differs from zero only in the neighbourhood
of one point (=c).

From (4) it appears that the coefficients A vanish, and from (5) that

Calling jpn dz, the whole momentum of the blow, Y, we

have (6),

and for the final solution

\[\sin(Fem,t) +...  \sin(F Me).\]

(7). In adapting this result to the case of a rod free at z=, we may
replace

by

If the blow be applied at a node of one of the normal components, that
component is missing in the resulting motion. The present calculation is
but a particular case of the investigation of Article 101.

169. As another example we may take the case of a bar, which is
initially at rest but deflected from its natural position by a lateral
force acting at =c. Under these circumstances the coefficients B vanish,
and the others are given by (4), Article 168.

Now

in which the terms free from the integral sign are to be taken between
the limits; by the nature of the case y, satisfies the same terminal
conditions as does u r, and thus all these terms vanish at both limits.
If the external force initially applied to the element dx be Y dx, the
equation of equilibrium of the bar gives

b = Y (1),

and accordingly

If we now suppose that the initial displacement is due to a force
applied in the immediate neighbourhood of the point x=c, we have

and for the complete value of y at time t, (2).

In deriving the above expression we have not hitherto made any special
assumptions as to the conditions at the ends, but if we now confine
ourselves to the case of a bar which is clamped at =0 and free at z=7,
we may replace

If we suppose further that the force to which the initial deflection is
due acts at the end, so that c=, we get

(3)

When t=0, this equation must represent the initial displacement. In
cases of this kind a difficulty may present itself as to how it is
possible for the series, every term of which satisfies the condition
y'\,'' =0, to represent an initial displacement in which this condition
is violated. The fact is, that after triple differentiation with respect
to a, the series no longer converges for 2=1, and accordingly the value
of y'' is not to be arrived at by making the differentiations first and
summing the terms afterwards. The truth of this statement will be
apparent if we consider a point distant dl from the end, and replace

u" (b-dl) by wu'' ()-u (0) di,

in which w' (/) is equal to

m 4/l 4 u(l)

For the solution of the present problem by normal coordinates the reader
is referred to Article 101.

170. The forms of the normal functions in the various particular cases
are to be obtained by determining the ratios of the four constants in
the general solution of

du/dx = m/l u.

If for the sake of brevity x' be written for m x/l the solution may

be put into the form

u=A ( (ar) + (a') + B ( (x') - (x'))

+ C ( (2) + (2')) + D ( (2) - (a)).

the hyperbolic cosine, \(\cosh(.)\), and sine \(\sinh(.)\)
are defined by the equations

\[\cosh z=( +e), \sinh(c)= e (2).\]

I have followed the usual notation, though the introduction of a special
symbol might very well be dispensed with, since

\[\cosh(x)  =\cos(i x), \sinh(x) =- i \sin(ix),\]
(3) where \(i=\sqrt{-1}\); and then the connection between the formulae of
circular and hyperbolic trigonometry would be more apparent. The rules
for differentiation are expressed in the equations

\[\cosh(x) = \sinh(z), a \sinh(.)  = \cosh(.) \]
\[\cosh() = \cosh(2), c \sinh(w) = \sinh(a) .\]

In differentiating (1) any number of times, the same four compound
functions as there occur are continually reproduced. The only one of
them which does not vanish with x' is (x) + (x), whose value is then 2.

Let us take first the case in which both ends are free. Since

Gu da and ae vanish with x, it follows that B= 0, D=0, so that

u=A ( (a) + (2') + C () (a') (4).

We have still to satisfy the necessary conditions when x=l, or x' =m.
These give
\[ A (- (m + (m) + C (- (m + \sinh(m)) = 0 . A( (m+ \sinh(m))+ C \cos(m) \]
(5),
equations whose compatibility requires that
\[ \cosh( m) - (m)2 = \sinh( 2 m) - (2) m, \]
or in virtue of the relation
\[ \cosh(2 m) - \sinh( 2 m) = 1 \]
(6),
m m = 1
(7).

This is the equation whose roots are the admissible values of m. If (7)
be satisfied, the two ratios of A: C given in (5) are equal, and either
of them may be substituted in (4). The constant multiplier being
omitted, we have for the normal function
\[ u = ( (m) - \sinh(m) () (a) \]
\[ - ( (m) - (m) (- + \sinh(8), or, if we prefer it, u = ( (m) - (m) eos \]
\[ + (a) + ( (m) + \sinh(m) (- + \sinh(9); \]
and the simple harmonic component of this type is expressed by

y = P u (F mt + c) (10).

a velocity depending only on the material of which the bar is formed,
and m is an abstract number. Hence for a given material and mode of
vibration the frequency varies directly as the radius of gyration of
the section about an axis perpendicular to the

171. The frequency of the vibration is m', in which 6 is plane of
bending, and inversely as the square of the length. These results might
have been anticipated by the argument from dimensions, if it were
considered that the frequency is necessarily determined by the value of
J, together with that of b-the only quantity depending on space, time
and mass, which occurs in the differential equation. If everything
concerning a bar be given, except its absolute magnitude, the frequency
varies inversely as the linear dimension.

These laws find an important application in the case of tuning forks,
whose prongs vibrate as rods, fixed at the ends where they join the
stalk, and free at the other ends. Thus the period of vibration of forks
of the same material and shape varies as the linear dimension. The
period will be approximately independent of the thickness perpendicular
to the plane of bending, but will vary inversely with the thickness in
the plane of bending. When the thickness is given, the period is as the
square of the length.

In order to lower the pitch of a fork we may, for temporary purposes,
load the ends of the prongs with soft wax, or file away the metal near
the base, thereby weakening the spring. To raise the pitch, the ends of
the prongs, which act by inertia, may be filed.

The value of 6 attains its maximum in the case of steel, for which it
amounts to about 5237 metres per second. For brass the velocity would be
less in about the ratio 15:1, so that a tuning fork made of brass would
be about a fifth lower in pitch than if the material were steel.

[For the design of steel vibrators and for rough determinations
of frequency, especially when below the limit of hearing, the
theoretical formula is often convenient. If the section of the bar
be rectangular and of thickness t in the plane of vibration,
and then with the above value of 6, and the values of m given
later, w:e get as applicable to the gravest mode

(clamped-free) frequency 84590
(free-free) frequency 538400 t P,

I and t being expressed in centimetres

The first of these may be used to ca late the pitch of steel
tuning-fork


The lateral vibrations of a bar may be excited by a blow,
when a tuning-fork is struck against a pad. This method is als'
employ for the harmonium, in which strips of metal or glass are
support at the nodes, in such a manner that the free vibrations
are but little impeded. A frictional maintenance may be obtained
with a bow, or by the action of the wetted finger upon a slender
rod of glass suitably attached. The electro-magnetic maintenance
of forks has been already considered, § 64. It may be applied with
equal facility to the case of metal bars, or even to that of
wooden planks carrying iron armatures, free at both ends and
supported at the nodes. The maintenance by a stream of wind
of the vibrations of harmonium and organ reeds may also be
referred to

The sound of a bar vibrating laterally may be reinforced by a
suitably tuned resonator, which may be placed under the middle
portion or under one end. On this principle dinner gongs have
been constructed, embracing one octave or more of the diatonic
scale.]


172. The solution for the case when both ends are clamped may be
immediately derived from the preceding by a double differentiation.
Since y satisfies at both ends the terminal conditions 2 2

it is clear that y'' satisfies dy''

y" =0, Ia = 9

which are the conditions for a clamped end. Moreover the general
differential equation is also satisfied by y''. Thus we may take,
omitting a constant multiplier, as before, u = ( (m) - (m) ) (2) - (2) -
( (m) - (m) (a) - (2) (1),


-while m is given by the same equation as before, namely, (m) (M = 1
(Q)e

We conclude that the component tones have the same pitch in the two
cases,

In each case there are four systems of points determined by the
evanescence of y and its derivatives. When y vanishes, there is a node;
where y/ vanishes, a loop, or place of maximum displacement; where y''
vanishes, a point of inflection; and where y vanishes, a place of
maximum curvature. Where there are in the first case (free-free) points
of inflection and of maximum curvature, there are in the second
(clamped-clamped) nodes and loops respectively; and vice versa, points
of inflection and of maximum curvature for a doubly-clamped rod
correspond to nodes and loops of a rod whose ends are free.

173. We will now consider the vibrations of a rod clamped at x=0, and
free at a=/. Reverting to the general integral (1) Article 170 we see
that A and C vanish in virtue of the conditions at = 0, so that

\[ u = B ( (a) - (a') + D ( (w - \sinh(w') (1). The remaining \]
conditions at 2=1 give
\[ B( (m-+ (m) + D ( (m) + \sinh(m) = 0 i B( - (m + \]
\[ \sinh(m) + D ( (m+ (m) = 0 \]
whence, omitting the constant multiplier,

\[ = ( (m) + \sinh(m) eos 7 - (a) \]

\[ - ( (m) + (m) (M2 \sinh(2), \]

or

wu = ( (m) + (m) eos - (ml)

\[ + ( (m) - \sinh(m) rin - (Ph .. (3); where m must be a root of (m) (m) + 1=0 \]
(4).

The periods of the component tones in the present problem are thus
different from, though, as we shall see presently, nearly related to,
those of a rod both whose ends are clamped, or free.

If the value of u in (2) or (3) be differentiated twice, the result
(u'\,') satisfies of course the fundamental differential equation.

a ow d aw'' . - a'' d ow bd

At 2=0, de" ae vanish, but ata=1 wand dz vanish. The function w'\,' is
therefore applicable to a rod clamped at 2 and free at 0, proving that
the points of inflection and of maximum curvature in the original curve
are at the same distances from the clamped end, as the nodes and loops
respectively are from the free end.

174. In default of tables of the hyperbolic cosine or its logarithm, the
admissible values of m may be calculated as follows. Taking first the
equation

(m) (m) = 1,

we see that m, when large, must approximate in value to 4(28 +1) 7, c
being an integer. If we assume

m= (264.1) we - (UB (2), 8 will be positive and comparatively small in
magnitude. - Substituting in (1), we find
\( \cot(8)=e"=e AGE De O-\)
or, if e" be called a,
\[ \arctan = 0P\]
(3), an equation which may be solved by successive approximation after
expanding \(\tan(a)\)  nd \(e''\) in ascending powers of the small quantity 8.
The result is
(4),
which is sufficiently accurate, even when 7=1. By calculation

B, = :0179666 - 0003228 + - 0000082 - 0000002 = 0176518.

B,, B,, B,, 8, are found still more easily. After A, the first term of
the series gives 8 correctly as far as six significant figures. The
table contains the value of , the angle whose circular measure is \(\beta\),
and the value of (/ 2), which will be required further on.

Dangling\footnote{ This process is somewhat similar to that adopted by Strehlke.}

Free-Free Bar.

1 10 -1 176518 1 0' 407-94 10 -2 88258 2 10 -3 777010 2' 40''-2699 10 -3
38850 3 10 -4 335505 6''-92029 10 -4 16775 4 10 -5 144989 299062 10 -6
72494 5 10 -7 626556 0129237 10 -7 31328

The values of m which satisfy (1) are

m1 = 4.7123890 + , = 4.7300408 m2 = 7.8539816 - , = 7.8532046 m3 =
10.9955743 + , = 10.9956078 m4 = 14.1371669 - , = 14.1371655 m5 =
17.2787596 + , = 17.2787596

after which m= (2 i +1) pi/2 to seven decimal places.

We will now consider the roots of the equation

(m) (m) =-1 (5).

Assuming

m= (2c-1) w- (HI) fa... (6)

we obtain the same result as before
\[ e m = \cot(1/2) = a' e^{} \]
where however \( a = e^{}\).

From this it appears that the series of values of is the same as that of
, though the corresponding suffixes are not the same. In fact

a=B,, 4=P,=

so that we have nothing further to calculate than , for which however
the series (4) is not sufficiently convergent
\footnote{This connexion between a and 8 does not appear to have been hitherto noticed.}.
The value of , may be obtained by trial and error from the equation

log,, cot 4, - 0.6821882 - 0.43429448 a 1 =0,

and will be found to be

= 0.3043077.

Another method by which m, may be obtained directly will be given
presently.

The values of m, which satisfy (5), are 15707963 + a, = 1875104

m, = 47123890 -a,= 4694737

m,= 78539816 + a,= 7854758

m, = 109955743 - a, = 10995541

m, = 141371669 + a, = 14137168 m, = 172787596 - a, = 17278759,

after which m = 4(2t- 1) sensibly.

The frequencies are proportional to m', and are therefore for the higher
tones nearly in the ratio of the squares of the odd numbers. However, in
the case of overtones of very high order, the pitch may be slightly
disturbed by the rotatory inertia, whose effect is here neglected.

175. Since the component vibrations of a system, not subject to
dissipation, are necessarily of the harmonic type, all the values of m ,
which satisfy

must be real. We see further that, if m be a root, so are also -m, mJ-1,
-m,/-1. Hence, taking first the lower sign, we have

n 5 (c08 m (m) = 1 m m = c - mn) ( m=) KC. (2).

If we take the logarithms of both sides, expand, and equate
coefficients, we get

eg) Zap Tae a

This is for a clamped-free rod.


From the known value of sum m, the value of m, may be derived

with the aid of approximate values of m,m, We find

sum m = -006547621,

and m, =000004237

m, = 000000069

m,= 000000005,

whence m, = 006543310

giving m, =1875105, as before.

In like manner, if both ends of the bar be clamped or free,

1- m 4/12.35 + ... = (1-m/m) (1- m/m) etc (4),

whence sum m4 = 1 = - etc, where of course the summation is exclusive of
the zero value of m.

176. The frequencies of the series of tones are proportional to m 2. The
interval between any tone and the gravest of the series may conveniently
be expressed in octaves and fractions of an octave. This is effected by
dividing the difference of the logarithms of m by the logarithm of 2.
The results are as follows :

1:4629 26478 24358 41332

31590 51036

37382, etc. 58288, etc.

where the first column relates to the tones of a rod both whose ends.are
clamped, or free ; and the second column to the case of a rod clamped at
one end but free at the other. Thus from the second column we find that
the first overtone is 26478 octaves higher than the gravest tone. The
fractional part may be reduced to mean semitones by multiplication by
12. The interval is then two octaves+ 77736 mean semitones. It will be
seen that the rise of pitch is much more rapid than in the case of
strings.

If a rod be clamped at one end and free at the other, the pitch of the
gravest tone is 2 (log 47300 - log 18751) + log 2 or 26698 octaves lower
than if both ends were clamped, or both free.

177. In order to examine more closely the curve in which the rod
vibrates, we will transform the expression for u into a form more
convenient for numerical calculation, taking first the case when both
ends are free. Since m= (2i+1)ar--(-1)8, (m= (f8, (m)= (imx (; and
therefore, m being a root of (m (m) (m) = (f).

Also

\[ \sinh(m) = m - 1 = \tan(m) = \cot(B) \]
or, since \( \cot(8)\)  is positive,
\[ \sinh(m) = \cot(f).  \]

Thus
\[ (m - \sinh(m) 1 - (tr) (8) (m) - (m (8) \]

( (8)- (tz (48) ( (8)- (im (8) ( (f + (tm (f)

(tr - (8) (8)(tr + (48)

We may therefore take, omitting the constant multiplier,

u = (8)(tr) + (8) ein - (wt)

- (8)(ta) - (8) () + (I)

= (m) (7)

+ (8) - (m) (dB) (1).

If we further throw out the factor ./2, and put /=1, we may take w= Fi
+F,+F, where F,= (tm ( ma -4ar+4(-1)8 log F,= ma loge + log (t8-log,/2
(2), log + F,=- mz log e + log (8) - log /2

from which u may be calculated for different values of c and x.


At the centre of the bar, =4, and Ff, F, are numerically equal in virtue
of e =cot 48. When 7 is even, these terms cancel. For F,, we have F, =
(-1) sin , which is equal to zero when is even, and to + 1 when 7 is
odd. When c is even, therefore, the sum of the three terms vanishes, and
there is accordingly a node in the middle.

When z=0, u reduces to - 2 (-1) (4 - 1) , which (since 8 is always small)
shews that for no value of 7 is there a node at the end. If a
long bar of steel (held, for example, at the centre) be gently tapped
with a hammer while varying points of its length are damped with the
fingers, an unusual deadness in the sound will be noticed, as the end is
closely approached.

178. Free-Free bar with three nodes.

Vibration with two nodes, i=1.,

If c=1, the vibration is the gravest of which the rod is capable. Our
formula become

=- ( x (270 + 1 0' 40'' 94) - 45 - 30' 2047
\(
log F,= 2054281 +37952391
\)
\(
log F, = - 2054231 a + 1:8494681,
\)
from which is calculated the following
table, giving the values of w for x equal to 00, 05, 10, etc.

The values of u:u(5) for the intermediate values of 2 (in the last column)
were found by interpolation formulae. If 0, p, g,7, 8,t be six
consecutive terms, that intermediate between g and r is x F, u:u(5)

000 + 7133200 + -0062408 + 7070793 + 1-4266401 + 1:645219

025 1:454176

050 -5292548 -0079059) -5581572; 1-0953179 1:263134

075 one 1072162

100 -3157243) -0100153 -4406005 7663401 8837528 125 6969004 "150 +
0846166) -0126874; -3478031 4451071 5133028 175 3341625 200 --1512020
0160726 -2745503 + :1394209 )+ -1607819 225 0054711 250 -+3786027
-0203609 -2167256 - -1415162 1631982 275 3109982 300 5849255 -0257934
-1710798 3880523 4475066 325 6714137 350 -7586838 -0326753 -1350477
5909608 6815032 375 1766629 400 8902038 -0413934 :1066045 7422059
8559210 425 9184491 450 +-9721635 -0524376 -0841519 8355740 9635940 475
: 9908730

500 - 1-000000 + 0664285 -0664282/- -867 1483 - 0000

Since the vibration curve is symmetrical with respect to the middle of
the rod, it is unnecessary to continue the table beyond x=5. The curve
itself is shewn in fig. 28.

Fig. 28.

To find the position of the node, we have by interpolation

ong , 1607819


which is the fraction of the whole length by which the node is distant
from the nearer end.

Vibration with three nodes. c= 2.

F, = ( (450 - 2' 40'' -27) - 45 + 1' 20185 log F,= 3410604 a + 44388816
log (- F,) = - 8'410604 + 18494850,

w:--u(0)

000 - 1-0000 025 8040 050 6079 075 4147 100 2274 125 - 0487 150 + 1175
175 2672 200 +3972 5037

In this table, as in the preceding, the values of u were calculated
directly for x = 000, 050, 100 etc., and interpolated for the
intermediate values. For the position of the node the table gives by
ordinary interpolation x=132. Calculating from the above formulae, we
find

u (1321) = - 000076,

u (1822) = + 000881, whence x = 132108, agreeing with the result
obtained by Strehlke. The place of maximum excursion may be found from
the derived function, We get

u' (3083) = + 0006077, uw' (3084) = - 0002227, whence u' (308373) = 0.
Hence u is a maximum, when = 308378; it then attains

the value 6636, which, it should be observed, is much less than the
excursion at the end.

The curve is shewn in fig. 29.

Fig. 29.

Vibration with four nodes. 7=3. F, = - ( (630 + 692) - 45 - 3-46 , log
F,= 4775832 x + 50741527, log F, = - 4775332 + 18494850. From this u (0)
=1:41424, u(.)=1:00579. The positions of the nodes are readily found by
trial and error. Thus u (3558) =--000087 su (3559) = + 001047,

whence is 0944, (Seebeck).

The position of the loop is best found from the derived

function. It appears that w'=0, when 2=-2200, and then

=- 9349. There is also a loop at the centre, where however the excursion
is not so great as at the two others.

Fig. 30.

a

We saw that at the centre of the bar F, and F, are numerically equal. In
the neighbourhood of the middle, F, is evidently very small, if c be
moderately great, and thus the equation for the nodes reduces
approximately to

ma T -F+(-) 6 + nm,

nm being an integer. If we transform the origin to the centre of the
rod, and replace m by its approximate value 4 (21+ 1) 7, we find

w +2n-7 tl 2vf4+1'
shewing that near the middle of the bar the nodes are uniformly spaced,
the interval between consecutive nodes being 2/ + (21+1). This
theoretical result has been verified by the measurements of Strehlke and
Lissajous.


For methods of approximation applicable to the nodes near the ends, when
c is greater than 3, the reader is referred to the memoir by Seebeck
already mentioned Article 160, and to Donkin's Acoustics (p. 194).

179. The calculations are very similar for the case of a bar clamped at
one end and free at the other. If wa F, and F=F,+F,+F,, we have in
general

F,= ( mx + ta +4(- ya ,

-1) (a) e; F,=-- (a) e,

V2 3 J2

If c = 1, we obtain for the calculation of the gravest vibration curve

ra

F, = (fe ma + 45 - 8 43-0666 ,
\[  log (- F,)= ma log e + 10300909, log (F) =- ma log e + 18444383.  \]
These give on calculation

F (0) = 000000, F( 6) = 743452, F (2) ="102974, F( 8) =1:169632, F (-4)
= 370625, F (1-0) = 1612224,
from which fig. 31 was constructed.

Fig. 31.

The distances of the nodes from the free end in the case of a

rod clamped at the other end are given by Seebeck and by Donkin.

2nd tone 0.2261.

3rd tone 0.1321, 0.4999.

4th tone 0.0944, 0.3558, 0.6439.

i-th tone

"The last row in this table must be understood as meaning that
(4j-3)/(4i-2) may be taken as the distance of the \(j\)th node from the free
end, except for the first three and the last two nodes.'' When both ends
are free, the distances of the nodes from the nearer end are 1st tone
0.2242.

2 tone 1321 5.

3 tone 0944 3558.

i-th tone 1.3222 4.9820 9.0007 4j-3 4i+2 4i+2 4i+2 4i+2

The points of inflection for a free-free rod (corresponding to the nodes
of a clamped-clamped rod) are also given by Seebeck;

1st point. 2nd point. th point.

No inflection point. 6000 "3593 5-0175 44+ 2

Except in the case of the extreme nodes (which have no corresponding
inflection-point, the nodes and inflection-points always occur in close
proximity.

180. The case where one end of a rod is free and the other supported
does not need an independent investigation, as it may be
referred to that of a rod with both ends free vibrating in an even mode,
that is, with a node in the middle. For at the central node y and y''
vanish, which are precisely the conditions for a supported end. In like
manner the vibrations of a clamped-supported rod are the same as those
of one-half of a rod both whose ends are clamped, vibrating with a
central node.

181. The last of the six combinations of terminal conditions occurs when
both ends are supported. Referring to (1) Article 170, we see that the
conditions at =0, give A =0, B=0; 80 that

\[ u= (C+D) (x +(C-D) \sinh(2),\]
Since u and w'' vanish when 2 = m, C- D=0,
and (m)=0. Hence the solution is

, a y = () - me Eb (1),

where 7 is an integer. An arbitrary constant multiplier may of course be
prefixed, and a constant may be added to t.

It appears that the normal curves are the same as in the case of a
string stretched between two fixed points, but the sequence of tones is
altogether different, the frequency varying as the square of 4 The nodes
and inflection-points coincide, and the loops (which are also the points
of maximum curvature) bisect the distances between the nodes.

182. The theory of a vibrating rod may be applied to illustrate the
general principle that the natural periods of a system fulfil the
maximum-minimum condition, and that the greatest of the natural periods
exceeds any that can be obtained by a variation of type. Suppose that
the vibration curve of a clamped-free rod is that in which the rod would
dispose itself if deflected by a force applied at its free extremity.
The equation of the curve may be taken to be y=- 81x + a',

which satisfies 4 = 0 throughout, and makes y and y' vanish at

0, and y'' at 7. Thus, if the configuration of the rod at time c be Y =
(- Bla + 2) (Mt (1),

the potential energy is by (1) Article 161, 6 (pt), while the
kinetic energy is 70 pol' p (pt; and thus p = T

Now p, (the true value of p for the gravest tone) is equal to

a x (1:8731) so that

shewing that the real pitch of the gravest tone is rather (but
comparatively little) lower than that calculated from the hypothetical
type. It is to be observed that the hypothetical type in question
violates the terminal condition y=0. This circumstance, however, does
not interfere with the application of the principle, for the assumed
type may be any which would be admissible as an initial configuration ;
but it tends to prevent a very close agreement of periods.

We may expect a better approximation, if we found our calculation on the
curve in which the rod would be deflected by a force acting at some
little distance from the free end, between which and the point of action
of the force (=c) the rod would be straight, and therefore without
potential energy. Thus potential energy = 6

The kinetic energy can be readily found by integration from the value of
y.

From 0 to c y=- 3cx + 2; and from c to / y=c'(c- 32),
as may be seen from the consideration that y and y' must not suddenly
change at z=c. The result is
kinetic energy = pw p (pt)E +4 (l-od (e+ 3r) , whence , po 6e8 E +s
(l-c) 30 be (2). The maximum value of 1+ p will occur when the point of
application of the force is in the neighbourhood of the node of the
second normal component vibration. If we take c= 31, we obtain a result
which is too high in the musical scale by the interval expressed by the
ratio 1 : 9977, and is accordingly extremely near the truth. This
example may give an idea how nearly the period of a vibrating system may
be calculated by simple means without the solution of differential or
transcendental equations.

The type of vibration just considered would be that actually assumed by
a bar which is itself devoid of inertia, but carries a load M at its
free end, provided that the rotatory inertia of M could be neglected. We
should have, in fact,

\[V= \cos(pt),  T= \sin(pt).\]

(8) Even if the inertia of the bar be not altogether negligible in
comparison with IM, we may still take the same type as the basis of an
approximate calculation :

\[V = l \cos(pt),\]



\[T = (2m + 5 +35) p \sin(pt),\]

so that , v=

whence (4),

that is, Mis to be increased by about one quarter of the mass of the
rod. Since this result is accurate when M is infinite, and does not
differ much from the truth, even when M =0, it may be regarded as
generally applicable as an approximation. The error will always be on
the side of estimating the pitch too high.

183. But the neglect of the rotatory inertia of M could not be justified
under the ordinary conditions of experiment. It is as easy to imagine,
though not to construct, a case in which the inertia of translation
should be negligible in comparison with the inertia of
rotation, as the opposite extreme which has just been considered.

If both kinds of inertia in the mass M be included, even though that of
the bar be neglected altogether, the system possesses two distinct and
independent periods of vibration.

Let 2 and denote the values of y and ay at =. Then the equation of the
curve of the bar is

and

Va IE (Bet - Bel + PO c (1);

while for the kinetic energy (2),

if ' be the radius of gyration of M about an axis perpendicular to the
plane of vibration. The equations of motion are therefore

+ ae

Mi - (62 - 318) = 0

ae

Mc + 2, (- 8lz + 276) =0

whence, if z and vary as (pt), we find

p= - a 14254 et C(4),

corresponding to the two periods, which are always different.

If we neglect the rotatory inertia by putting x' =0, we fall back on our
previous result

The other value of p is then infinite.

If ' : 1 be merely small, so that its higher powers may be neglected,


If on the other hand be very great, so that rotation is prevented,
(6),
the latter of which is very small. It appears that when rotation is
prevented, the pitch is an octave higher than if there were no rotatory
inertia at all. These conclusions might also be derived
directly from the differential equations; for if ' = Ps ,9=0, and then
but if '=0, o=3, by the second of equations (3), and in that case


184. If any addition to a bar be made at the end, the period of
vibration is prolonged. If the end in question be free, suppose first
that the piece added is without inertia. Since there would be no
alteration in either the potential or kinetic energies, the pitch would
be unchanged; but in proportion as the additional part acquires inertia,
the pitch falls (Article 88).

In the same way a small continuation of a bar beyond a clamped end would
be without effect, as it would acquire no motion. No change will ensue
if the new end be also clamped ; but as the first clamping is relaxed,
the pitch falls, in consequence of the diminution in the potential
energy of a given deformation.

The case of a `supported' end is not quite so simple. Let the original
end of the rod be .A, and let the added piece which is at first supposed
to have no inertia, be 4B. Initially the end A is fixed, or held, if we
like so to regard it, by a spring of infinite stiffness. Suppose that
this spring, which has no inertia, is gradually relaxed. During this
process the motion of the new end B diminishes, and at a certain point
of relaxation, B comes to rest. During this process the pitch falls. B,
being now at rest, may be supposed to become fixed, and the abolition of
the spring at A entails another fall of pitch, to be further increased
as AB acquires inertia.

185. The case of a rod which is not quite uniform may be
treated by the general method of Article 90. We have in the notation
there adopted

a, = I pe ot, der, da, = I Spo u, dz,

whence, P, being the uncorrected value of p,

[If the motion be strictly periodic with respect to a, V' is
proportional to and both quantities vanish at a node.
Accordingly an irregularity situated at a node of this kind of motion
has no effect upon the period. A similar conclusion will hold good
approximately for the interior nodes of a bar vibrating with
numerous subdivisions, even though, as when the terminals are
clamped or free, the mode of motion be not strictly periodic with
respect to xJ]

% If the rod be clamped at 0 and free at J,
For example, if the rod be clamped at 0 and free at J,

The same formula applies to a doubly free bar. The effect of a small
load dM is thus given by B,m' udM tS ek: ) 2), P pe, I uj M ) where M'
denotes the mass of the whole bar. If the load be at
the end, its effect is the same as a lengthening of the bar in the ratio
M' : M'+ dit. (Compare Article 167.)

[In (2) dM is supposed to act by inertia only ; but a similar
formula may conveniently be employed when an irregularity of
mass dM depends upon a variation of section, without a change
of mechanical properties. Since B -

so that the effect of a local excrescence is given by

If the thickness in the plane of bending be constant,
and

Further,
and thus

If, however, the thickness in the plane perpendicular to that of
bending be constant, and in the plane of bending variable (27),
then 8 = 87 770 = 887 70 = 3 8 a)o;
and in place of (4)
+ W
If a tuning-fork be filed (dM negative) near the stalk (clamped end),
the pitch is lowered ; and if it be filed near the free end, the
pitch is raised. Since the effects of a given stroke of

the file are equal and opposite in the circumstances of (4), but in
the circumstances of (5) the effect at the stalk is three times as
great as at the free end ]

186. Correction for Rotatory Inertia.
The same principle may be applied to estimate the correction due to
the rotatory inertia of a uniform rod. We have only to find what
addition to make to the kinetic energy, supposing that the bar vibrates
according to the same law as would obtain, were there no rotatory
inertia.

Let us take, for example, the case of a bar clamped at 0 and free at J,
and assume that the vibration is of the type, y = os pt,

where wu is one of the functions investigated in Article 179. The
kinetic energy of the rotation is

= ip" (pt)(2u u' + ''),

by (2) Article 165.

To this must be added U e p
(pt) u' dx, or pet `p' (pt u; ; 0

so that the kinetic energy is increased in the ratio

1475 (2) +m).

The altered frequency bears to that calculated without allowance for
rotatory inertia a ratio which is the square root of the reciprocal of
the preceding. Thus (1). By use of the relations

\[\cosh(m) = - \sec(m), \sinh(m) = \cos(7) \tan(m),\]

we may express uv' : u when =/ in the form

\[wu - \sin(m) \cos() \]

if we substitute for m from m= (2t--1)r-(-1)'a In the case of the
gravest tone, a =3043, or, in degrees and minutes, a= 17 26', whence

Thus

which gives the correction for rotatory inertia in the case of the
gravest tone.

When the order of the tone is moderate, a is very small, and then uw
:u=1 sensibly,

and p:P=1 -(1+5) mel (3),

shewing that the correction increases in importance with the order of
the component.

In all ordinary bars : J is very small, and the term depending on its
square may be neglected without sensible error.

187. When the rigidity and density of a bar are variable from point to
point along it, the normal functions cannot in general be expressed
analytically, but their nature may be investigated by the methods of
Sturm and Liouville explained in Article 142.

If, as in Article 162, B denote the variable flexural rigidity at any
point of the bar, and pw dz the mass of the element, whose length is dx,
we find as the general differential equation

a ay ay

the effects of rotatory inertia being omitted. If we assume that
\(y  \cos(rt)\), we obtain as the equation to determine the form of the
normal functions

z (B 74) SVD Y v (2),

in which v' is limited by the terminal conditions to be one of an:
infinite series of definite quantities v,', v,v,

Let us suppose, for example, that the bar is clamped at both ends, so
that the terminal values of y and ay vanish. The first normal function,
for which v has its lowest value ,', has no internal root, so that the
vibration-curve lies entirely on one side of the equilibrium-position.
The second normal function has one internal root, the third function has
two internal roots, and, generally, the r function has r - 1 internal
roots.

Any two different normal functions are conjugate, that is to say, their
product will vanish when multiplied by pw dx, and integrated over the
length of the bar.

Let us examine the number of roots of a function f(x) of the form

F (2) =b

compounded of a finite number of normal functions, of which the function
of lowest order is u,(x) and that of highest order is u, (x). If the
number of internal roots of f (x) be \(\mu\), so that there are +4 roots
in all, the derived function:/'' (2) cannot have less than 4+ 1 internal
roots besides two roots at the extremities, and the second derived
function cannot have less than 4+ 2 roots

te


No roots can be lost when the latter function is multiplied by B, and
another double differentiation with respect to a will leave at least
internal roots. Hence by (2) and (3) we conclude that

Ve Pa Um (x) + Vest. Pats Une (z) + set v,. , Ue (z)...(4)

has at least as many roots as f(a''). Since (4) is a function of the
same form as f(a), the same argument may be repeated, and a series of
functions obtained, every member of which has at least as many roots as
f(x) has. When the operation by which (4) was derived from (3) has been
repeated sufficiently often, a function is arrived at whose form differs
as little as we please from that of the component normal function of
highest order u(.); and we conclude that f(z) cannot have more than n-1
internal roots. In like manner we may prove that f(a) cannot have less
than m-1 internal roots.

The application of this theorem to demonstrate the possibility of
expanding an arbitrary function in an infinite series of normal
functions would proceed exactly as in Article 142.

[An analytical investigation of certain cases where the section
of a rod is supposed to be variable, will be found in a memoir by
Kirchhoff ].

188. When the bar, whose lateral vibrations are to be considered, is
subject to longitudinal tension, the potential energy of any
configuration is composed of two parts, the first depending on the
stiffness by which the bending is directly opposed, and the second on
the reaction against the extension, which is a necessary accompaniment
of the bending, when the ends are nodes. The second part is similar to
the potential energy of a deflected string; the first is of the same
nature as that with which we have been occupied hitherto in this
Chapter, though it is not entirely independent of the permanent tension.

Consider the extension of a filament of the bar of section dw, whose
distance from the axis projected on the plane of vibration is . Since
the sections, which were normal to the axis originally, remain normal
during the bending, the length of the filament bears to the
corresponding element of the axis the ratio R+7: R, RB being the radius
of curvature. Now the axis itself is extended in the ratio g:qg+T,
reckoning from the unstretched state, if Tw denote the whole tension to
which the bar is subjected.

Hence the actual tension on the filament is T+h(P + ade) from which we
find for the moment of the couple acting across the section r + r+3 9 n
dy = 14 to,

and for the whole potential energy due to stiffness

4(q+T) ef (F2) di. (L),

an expression differing from that previously used (Article 162) by the
substitution of g+ 7' for q.

Since g is the tension required to stretch a bar of unit area to twice
its natural length, it is evident that in most Practical cases T would
be negligible in comparison with q.

The expression (1) denotes the work that would be gained during the
straightening of the bar, if the length of each element of the axis were
preserved constant during the process. But when a stretched bar or
string is allowed to pass from a displaced
to the natural position, the length of the axis is decreased. The
amount of the decrease is 4 (2) dx, and the corresponding gain d: + To
(2) ae. Thus ,

(2).

The variation of the first part due to a hypothetical displacement is
given in Article 162. For the second part, we have

12 (ff arn BRO Hi Boe

In all the cases that we have to consider, Sy vanishes at the limits.
The general differential equation is accordingly

a ay dy 2 q+ 2) oY 7 Th 4 Poa =p a = 0 or, if we put g+T=b'p. T=a"p, .
d'y dy dy 2(eg- az) - a 54+ Gi =0 (4).

For a more detailed investigation of this equation the reader is
referred to the writings of Clebsch\footnote{Theorie der Elasticitat fester
Korper. Leipzig, 1862} and Donkin.

189. If the ends of the rod, or wire, be clamped; dy/dx =0, and the
terminal conditions are satisfied. If the nature of the support be such
that, while the extremity is constrained to be a node, there is no
couple acting on the bar, d2 y/ dx2 must vanish, that is to say, the end
must be straight. This supposition is usually taken to represent the
case of a string stretched over bridges, as in many musical instruments;
but it is evident that the part beyond the bridge must partake of the
vibration, and that therefore its length cannot be altogether a matter
of indifference.

If in the general differential equation we take y proportional to (nt),
we get

=0 (1),

which is evidently satisfied by

y = (i x/l) (nt) (2),

if n be suitably determined. The same solution also makes y and y''
vanish at the extremities. By substitution we obtain for n,

n = (3)

which determines the frequency.

If we suppose the wire infinitely thin, n n =7 m'a'/(l2), the same as
was found in Chapter VI., by starting from the supposition of perfect
flexibility. If we treat \(\kappa :l\) as a very small quantity, the
approximate value of n is

n =

For a wire of circular section of radius r, \(\kappa 2 =(1/4) r2\), and
if we replace b and a by their values in terms of q. T, and \(\rho\),

n = (4),

which gives the correction for rigidity\footnote{Donkin's Acoustics, Art. 184}.
Since the expression within brackets involves c, it appears that the harmonic
relation of the component tones is disturbed by the stiffness.

190. Permanent tension.
The investigation of the correction for stiffness when the ends of
the wire are clamped is not so simple, in consequence of the change of
type which occurs near the ends. In order to pass from the case of the
preceding section to that now under consideration an additional
constraint must be introduced, with the effect of still further raising
the pitch. The following is, in the main, the investigation of Seebeck and Donkin.

If the rotatory inertia be neglected, the differential equation becomes

(Dt = Sip De ppg) Y= 0 (1),

where D stands for d/dx. In the equation

p- p-,=0, x'b = 0,

one of the values of D2 must be positive, and the other negative. We may
therefore take

D'- 5 Dt f= (Dt al) (D +B) (2),

and for the complete integral of (1)

\[ y=A (x) + B \sinh(x) + C (x) + D (x) \]
(3),

where and are functions of n determined by (2).

The solution must now be made to satisfy the four boundary conditions,
which, as there are only three disposable ratios, lead to an equation
connecting , beta, l. This may be put into the form

\[ \sinh(al) (Bl re j) - (al) (pl)\]
(4) The value of,
determined by (2), is
\[ \sinh(al) fl  - \cosh(al) \]
(5).  From (2) we find also that (6)

Thus far our equations are rigorous, or rather as rigorous as the
differential equation on which they are founded; but we shall now
introduce the supposition that the vibration considered is but slightly
affected by the existence of rigidity. This being the case, the
approximate expression for y is

y= = (22) (F at)

and therefore

= i / l, n = i / l (7).

nearly.

The introduction of these values into the second of equations (6) proves
that n Be or ae is a small quantity under the circumstances
contemplated, and therefore that a' is a large quantity. Since (l),
\( \sinh(l) \)
are both large, equation (5) reduces to
\[ \tan(l) = 2 n b \kappa /a2.  \]
or, on substitution of the approximate value for 8 derived from (6),
\[ \tan( n l/ a) = 2 n b \kappa/ a2 \]
The approximate value of n l/a is i . If we take
\(n l = i \pi  + \theta\) , we get

\[tan (i \pi + \theta ) = \tan(\theta ) = \theta =
 2 n b \kappa/ a2 = 2 i \pi (b/a) \kappa / l\]
so that,
\[ n = i a/l (1 + 2 (b/a) \kappa/l).  \]
(8) According to this equation the component tones are all raised in pitch
by the same small interval, and therefore the harmonic relation is not
disturbed by the rigidity. It would probably be otherwise if terms
involving x': were retained ; it does not therefore follow that the
harmonic relation is better preserved in spite of rigidity when the ends
are clamped than when they are free, but only that there is no
additional disturbance in the former case, though the absolute
alteration of pitch is much greater. It should be remarked that 6:a or
/(q+T7) : /T7, is a large quantity, and that, if our result is to be
correct, : 1 must be small enough to bear multiplication by 6 : a and
yet remain small.

The theoretical result embodied in (8) has been compared with experiment
by Seebeck, who found a satisfactory agreement. The constant of
stiffness was deduced from observations of the rapidity of the
vibrations of a small piece of the wire, when one end was clamped in a vice.


\todo[inline]{ paste from 2nd edition }

Dangling\footnote{As the result of a second approximation Seebeck gives (Zoc. cii.)}

191. Resultant of two trains of waves.
It has been shewn in this chapter that the theory of bars, even
when simplified to the utmost by the omission of unimportant quantities,
is decidedly more complicated than that of perfectly flexible strings.
The reason of the extreme simplicity of the vibrations of strings is to
a velocity independent of the wave length, so that an arbitrary wave is
allowed to travel without decomposition. But when we pass from strings
to bars, the constant in the differential equation, viz. oY 5 'b ay. 0,
is no longer expressible as a velocity, and therefore the velocity of
transmission of a train of harmonic waves cannot depend on the
differential equation alone, but must vary with the wave length. Indeed,
if it be admitted that the train of harmonic waves can be propagated at
all, this consideration is sufficient by itself to prove that the
velocity must vary inversely as the wave length. The same thing may be
seen from the solution applicable to waves propagated in one direction,
viz. y = (Vt- 2), which satisfies the differential equation if

V= 2 b/ (1).

Let us suppose that there are two trains of waves of equal amplitudes,
but of different wave lengths, travelling in the same direction. Thus

x


t t y = 00s 2x + (2a)

1 31 1 1 1,1 1 = 2 (me(; -=)- -2 (5 -=) (2+2)- a( + +). (2). If r' -7, '
-2 be small, we have a train of waves, whose amplitude slowly varies
from one point to another between the values 0 and 2, forming a series
of groups separated from one another by
regions comparatively free from disturbance. In the case of a string or
of a column of air, varies as t, and then the groups move
forward with the same velocity as the component trains, and there is no
change of type. It is otherwise when, as in the case of a bar vibrating
transversely, the velocity of propagation is a function of the wave
length. The position at time c of the middle of the group which was
initially at the origin is given by

1 1 1 1 tC -2)-(5-x) =
which shews that the velocity of the group is 1 1 ,/11).fl) ,,f (-3) G-x)=8G) 8G)
If we suppose that the velocity V of a train of waves varies
as MM", we find 1 4 a() 4G) 1, ,/l a(3) a(;) In the present case n =- 1,
and accordingly the velocity of the groups is twice that of the
component waves'.

(MAU) (8)e

192. On account of the dependence of the velocity of propagation on the
wave length, the condition of an infinite bar at any time subsequent to
an initial disturbance confined to a limited portion, will have none of
the simplicity which characterises the corresponding problem for a
string; but nevertheless Fourier's investigation of this question may
properly find a place here.

It is required to determine a function of and , so as to satisfy

dy dy ae + ad O (1), and make initially y= c (a), y= (a). A solution of
(1) is Y = (QE (G (L- A). (2), where g and a are constants, from which
we conclude that
y= da Fe) dg (qt (g (7 - a)

Dangling\footnote{ In the corresponding problem for waves on the surface of deep water,
the velocity of propagation varies directly as the square root of the
wave length, so that n=. The velocity of a group of such waves is
therefore one half of that of the component trains.}

\todo[inline]{ paste from 2nd edition }
footnote [See note on Progressive Waves, appended to this volume]

is also a solution, where F() is an arbitrary function of . If now we
put t = 0,

yo = dx F(q) dq (q (x - ),

which shews that F() must be taken to be (1/(2 ))( ), for then by
Fourier's double integral theorem yo =(x). Moreover, y=0; hence

1 f '' 2 y= da a) dq (q t (q ( -) + (3) satisfies the
differential equation and makes initially,

By Stokes' theorem (Article 95), or independently, we may now supply the
remaining part of the solution, which has to satisfy the differential
equation while it makes initially y=0, y= (a); it is

y= (4).

The final result is obtained by adding the right-hand members of (3) and
(4).

In (8) the integration with respect to q may be effected by means of the
formula

aq (g) (e) =4 (F) (5),
which may be proved as follows. If in the well-known integral formula

we put x+b for x, we get

Now suppose that a=t=e"", where \(i=\sqrt -1\), and retain only the real
part of the equation. Thus

whence

(at) (2h) = (),
from which (5) follows by a simple change of variable. Thus equation (3)
may be written
y=
or if (x-x)/(2 t) = ,
y= (6).

\section{Circular Rings} \label{sect:192a} 192 a.
If the axis of the rod be curved instead of straight, we obtain
problems which may be regarded as extensions of those of the present and
of the last chapters. The most important case under this head is
that of a circular ring, whose section we will regard as also circular
and of radius (c) small in comparison with the radius (a) of the
circular axis. The investigation of the flexural modes of vibration,
executed in the plane of the ring, is analogous to the case of a
cylinder (see § 233), and was first effected
by Hoppe \footnote{Crelle, Bd. 63, p. 158, 1871}. If s be the
number of periods in the circumference, the coefficient p of the time in
the expression for the vibrations is given by

. 4 l + pa (1),
where q is Young's modulus and p the density of the material.
This may be compared
with equation (9) § 233. To fall back upon the case of a straight axis
we have, only to suppose 8 and a to be infinite in such a manner that
\(2 \pi a/s\)
is equal to the proposed linear period. The vibrations in
question are then purely transverse. In the class of vibrations
considered above the circular axis remains unextended, and (§ 232) the
periods are comparatively long. For the other class of vibrations in the
plane of the ring, Hoppe found
(2).
The frequencies are here independent of c, and the vibrations are analogous
to the longitudinal vibrations of straight rods.

If s 0 in (2), we have the solution for vibrations which are purely
radial.

For flexural vibrations perpendicular to the plane of the ring, the
result \footnote{Michell, Messenger of Mathematics, XIX, 1889}
corresponding to (1) is (3), the difference consisting only in
the occurrence of Poisson's ratio
\(\mu \) in the denominator.

Our limits will not allow of our dwelling further upon the problem of
this section. A complete investigation will be found in Love's Treatise
on Elasticity, Chapter XVIII. The effect of a small curvature upon the
lateral vibrations of a limited bar has been especially considered by
Lamb \footnote{Proc. Lond. Math.  Soc., XIX., p. 365, 1888.}.
% next chapter is  CHAPTER IX. VIBRATIONS OF MEMBRANES.
